% 08-triaxial.tex -- the deformed source: spin and tides.
% Input by mostlyharmless.tex.  No preamble here.

\section{The deformed source: spin and tides}
\label{sec:triaxial}

Sections~\ref{sec:conventions}--\ref{sec:full-map} answer one question: given a surface
brightness map on the unit sphere, what is its complex visibility.  This section supplies
the map.  A star deformed by its own rotation and by a companion has a shape and a surface
brightness that are both consequences of a single potential, and the two are not
independent: on a level surface the normal lies along the gradient, so the surface gravity
follows from the shape.  Everything here is dimensionless in units where $GM = 1$ and the
undeformed radius is $1$.

\subsection{The quadrupole potential}
\label{sec:deformation-potential}

Write $n = \mathbf{x}/r$ for the direction, with $\mathbf{x}$ the position in space, and
$P_2(\xi) = \tfrac12(3\xi^2-1)$ for the Legendre polynomial.  Bold symbols are points in
space throughout; plain $x$ and $y$ are the sky coordinates of
section~\ref{sec:conventions}, so that $x_i$ and $y_i$ are the North and East components
of a sample's sky position and are not components of a position vector.  The
potential is a monopole plus two quadrupoles and one constant,
\begin{equation}
  \Phi(\mathbf{x}) \;=\; -\frac{1}{r}
    \;-\; \epsilon_\text{tide}\, r^2 P_2(n\cdot \hat e_\text{tide})
    \;+\; \epsilon_\text{spin}\, r^2 P_2(n\cdot \hat e_\text{spin})
    \;-\; \epsilon_\text{spin}\, r^2 .
  \label{eq:deformation-potential}
\end{equation}
Both perturbations are pure $l = 2$, and the amplitudes are positive for a tidal bulge
along $\hat e_\text{tide}$ and a centrifugal flattening about $\hat e_\text{spin}$.  The
last term is the $l = 0$ part of the centrifugal potential: the centrifugal term is not a
pure quadrupole, because the frame rotating with the star is not inertial, and its $l = 0$
piece renormalises the effective monopole.  Dropping it would leave the star's mean radius
wrong by $\epsilon_\text{spin}$.

Amplitudes are quoted as $A r_0^3/GM$, so $\epsilon = 0.2$ means a twenty per cent
deformation of the \emph{potential} at the reference radius, whatever axis ratio that
produces.  The surface is the level set $\Phi = -1$.  The volume-equivalent sphere is a
separate quantity, since the deformation changes the volume, and the two should not be
confused when quoting an axis ratio.

\subsection{The level surface, and when it closes}
\label{sec:deformation-closure}

Along a direction $n$ the level condition is one equation in one unknown,
\begin{equation}
  \frac{1}{r} + E(n)\, r^2 \;=\; 1 ,
  \label{eq:level-condition}
\end{equation}
where the depth of the perturbation there is
\begin{equation}
  E(n) \;=\; \epsilon_\text{tide} P_2(n\cdot\hat e_\text{tide})
           - \epsilon_\text{spin}\bigl[P_2(n\cdot\hat e_\text{spin}) - 1\bigr] .
  \label{eq:level-depth}
\end{equation}
The sign of the spin term is negative because its coefficient in
\eqref{eq:deformation-potential} is positive while it \emph{lowers} the potential.

The left side of \eqref{eq:level-condition} is not monotone in $r$: the monopole falls and
the quadrupole grows.  Its minimum is at $r^3 = 1/(2E)$ and equals
$1.89\,E^{1/3}$, so a solution exists only while
\begin{equation}
  E(n) \;<\; 0.148 \qquad \text{for every } n ,
  \qquad 0.148 = (1/1.89)^3 .
  \label{eq:closure}
\end{equation}
Past that the surface opens along the deepest direction, as a Roche lobe does.  This is a
\emph{physical} limit of the model, not a numerical one, and the amplitude that triggers it
is the useful result: in the perpendicular case, where the depth along the tide axis is
\begin{equation}
  E \;=\; \epsilon_\text{tide} + 1.5\,\epsilon_\text{spin} ,
  \label{eq:equatorial-depth}
\end{equation}
the bound reads $0.148$ for the tide alone, $0.099$ for the spin alone (whose coefficient
is $1.5\epsilon_\text{spin}$), and $0.059$ for two equal amplitudes.  A twenty per cent
deformation of the potential, which was the first target, is therefore not attainable in
this model at all.

When the two axes are not perpendicular but separated by an angle $\theta$, the deepest
direction is found from the quadratic form \eqref{eq:depth-form} below, and the depth
becomes
\begin{equation}
  E_\text{max} \;=\; \tfrac34\Bigl[(\epsilon_t-\epsilon_s)
    + \sqrt{(\epsilon_t-\epsilon_s)^2 + 4\epsilon_t\epsilon_s\cos^2\theta}\Bigr]
    + \tfrac{3\epsilon_s-\epsilon_t}{2} ,
  \label{eq:depth-latitude}
\end{equation}
which reduces to \eqref{eq:equatorial-depth} at $\theta = \pi/2$ and to
$\epsilon_\text{tide}$ alone at $\theta = 0$.  A tide aligned with the spin costs nothing in
the centrifugal term, which has no gradient along the spin axis, so tilting the tide
\emph{loosens} the bound.  That is the opposite of the intuition that a more oblique system
is more constrained, and it is what makes the latitude a useful parameter rather than a
nuisance.

\subsection{The surface gravity}
\label{sec:deformation-gravity}

The gradient is worth writing out explicitly.  The $\Phi$ of
\eqref{eq:deformation-potential} simplifies once $r^2(n\cdot\hat e)^2 =
(\mathbf{x}\cdot\hat e)^2$ is used to move the $r$-dependence out of the Legendre
polynomials:
\begin{equation}
  \Phi(\mathbf{x}) \;=\; -\frac{1}{r}
    \;+\; \tfrac12(\epsilon_\text{tide} - 3\epsilon_\text{spin})\,r^2
    \;-\; \tfrac32\epsilon_\text{tide}\,(\mathbf{x}\cdot\hat e_\text{tide})^2
    \;+\; \tfrac32\epsilon_\text{spin}\,(\mathbf{x}\cdot\hat e_\text{spin})^2 ,
  \label{eq:deformation-potential-quadratic}
\end{equation}
which differentiates term by term to
\begin{equation}
  \nabla\Phi \;=\; \frac{\mathbf{x}}{r^3}
    \;+\; (\epsilon_\text{tide} - 3\epsilon_\text{spin})\,\mathbf{x}
    \;-\; 3\epsilon_\text{tide}\,(\mathbf{x}\cdot\hat e_\text{tide})\,\hat e_\text{tide}
    \;+\; 3\epsilon_\text{spin}\,(\mathbf{x}\cdot\hat e_\text{spin})\,\hat e_\text{spin} .
  \label{eq:deformation-gradient}
\end{equation}
Evaluating \eqref{eq:deformation-gradient} at $\mathbf{x} = r(n)\,n$, with $r(n)$ from
\eqref{eq:level-condition}, and taking its norm gives $g(n)$.  The direction of
$\nabla\Phi$ is of course the outward normal used for the visibility sign in
section~\ref{sec:deformation-exact}, so the same expression serves both purposes.

Two checks.  With both amplitudes zero the first term survives and $|\nabla\Phi| = 1/r^2$,
which is $1$ at $r = 1$ as required.  And to first order the perturbation reaches $g$ by two
routes.  Write $\phi_1$ for the perturbation to the potential, $\phi_1(r,n) = -E(n)\,r^2$ with
$E$ the depth of \eqref{eq:level-depth} --- the sign being the one that makes $1/r + E r^2$ the
level condition --- so that on the reference sphere $\phi_1 = -E(n)$, and $g = 1 + 4\phi_1$.
Holding
$r$ fixed, the amplitude's own radial derivative contributes
$\partial_r\phi_1 = 2\phi_1$, since $\phi_1 \propto r^2$.  And the surface itself moves,
$\delta r = -\phi_1$, which the monopole's $1/r^2$ turns into a further $-2\delta r = 2\phi_1$.
The two are equal, so
\begin{equation}
  \delta g \;=\; 2\phi_1 - 2\,\delta r \;=\; 4\phi_1 ,
  \label{eq:deformation-gravity-shift}
\end{equation}
the four being a quadrupole signature rather than a coincidence: for a perturbation
$\phi_1 \propto r^n$ the same argument gives $\delta g = (n+2)\phi_1$, so an octupole would
give five.  Writing $\phi_1$ out,
\begin{equation}
  g \;=\; 1 - 4\epsilon_\text{tide}\,P_2(n\cdot\hat e_\text{tide})
    + 4\epsilon_\text{spin}\bigl[P_2(n\cdot\hat e_\text{spin}) - 1\bigr]
    + O(\epsilon^2) ,
  \label{eq:deformation-gravity-first-order}
\end{equation}
where $\epsilon$ stands for both amplitudes jointly, so the remainder contains the squares and
the cross term.  This is the same order as the rest of the deformation model ---
$\delta r = -\phi_1$ is linear too --- so it introduces no inconsistency.  It is a check on the exact gradient rather than a substitute
for it: the code evaluates $|\nabla\Phi|$ from \eqref{eq:deformation-gradient}, and this is
what certificate \vtag{19} compares it against.  The sign is the physically expected one: a
rotating star is brightest at its poles and dimmest at its equator, six
$\epsilon_\text{spin}$ below the undeformed value there.
\subsection{Surface brightness}
\label{sec:deformation-brightness}

The surface gravity is $g = |\nabla\Phi|$ evaluated on the level surface, and it equals
$1$ for the undeformed star, so the brightness falls out dimensionless:
\begin{equation}
  I \;\propto\; g^{p} , \qquad g = |\nabla\Phi| .
  \label{eq:brightness-law}
\end{equation}

The bolometric von Zeipel law is $p = 1$.  Physically motivated laws have $p \le 1$: a
convective envelope is nearer to $p \approx 0.3$, because the flux there is carried by
convection and so responds to gravity more weakly.  Thus $p = 1$ is an upper bound rather
than a convenient default, and it is the one the code uses.

Limb darkening is \emph{absent} by construction: \eqref{eq:brightness-law} depends only on
gravity, not on the angle between the surface normal and the line of sight.  That is a
deliberate scope restriction until the limb-darkening stage, and any figure caption or
comparison with a resolved observation has to say so.


Limb darkening multiplies this by a function of the viewing angle, and the code applies it
in the same place to every row: the sampled brightness, the linearised pattern and the fitted
coefficients all carry the factor
\begin{equation}
  1 - u_\text{limb}\,(1-\mu) , \qquad \mu = \hat n\cdot\hat e_\text{obs} ,
  \label{eq:limb-factor}
\end{equation}
so $u_\text{limb} = 0$ recovers \eqref{eq:brightness-law} and $u_\text{limb} = 1$ darkens the
limb to nothing.  It is linear in $\mu$ rather than a power law because $\mu$ is a single
harmonic, which is what keeps the coefficients of section~\ref{sec:model-map} closed; the
first-order coefficients of \eqref{eq:limb-coefficients} treat it as a dipole about the
observer's axis.

\paragraph{Why the limb is not done analytically.}  The original plan was to expand the star
and the limb-darkening factor separately and multiply them with Gaunt coefficients, which
would be fully analytic.  On the disk that product is cheaper than it looks: for a disk seen
face-on
\begin{equation}
  \mu \;=\; \sqrt{1-\rho^2}
  \label{eq:limb-mu-disk}
\end{equation}
is a function of the impact parameter alone, so a limb law is a radial profile, and an
expansion in Zernike modes with $m = 0$ needs no Gaunt at all --- $\mu$ is already the
variable the basis is built from.  That is why the classical visibility of a linearly
limb-darkened disk is a two-term combination of Bessel functions.  On the sphere $\mu$ is
instead a degree-one factor, and re-expanding its product with the map's harmonics needs the
$2\times1 \to 1$ and $2\times1 \to 3$ coefficients, which is where the Gaunt coefficients would
enter.

What settled the question was not the limb but the shape.  Spin and tide make the silhouette
and the brightness both depend on $\epsilon$, with no closed form for the level surface, so
the analytic route has no route through the shape.  The row that is implemented therefore fits
the rendered image instead, and carries the limb as a factor rather than as a product of
expansions.

\subsection{The exact visibility}
\label{sec:deformation-exact}

With the surface and the brightness fixed, the visibility is the sum over the sampled
surface that section~\ref{sec:mode-visibility} defines, weighted by the sky-plane area each
sample projects to:
\begin{equation}
  V(u,v) \;=\; \sum_i I(n_i)\,\frac{(N_i\cdot\hat e_\text{obs})\, r_i^2}{n_i\cdot N_i}\,
    \Delta\Omega_i \; e^{-2\pi i (u x_i + v y_i)} .
  \label{eq:deformation-visibility}
\end{equation}
Two properties of a level surface make this cheap.  The outward normal $N$ \emph{is} the
gradient direction, so visibility is the sign of $N\cdot\hat e_\text{obs}$: there is no
z-buffer and no nested root find.  And for modest deformation the body is convex, so the
map from direction to sky position is one-to-one and the near hemisphere is simply the set of
directions with $N\cdot\hat e_\text{obs} > 0$.

Equation~\eqref{eq:deformation-visibility} is a quadrature, not a periodic transform, so its
error is ordinary discretisation error that falls with the sampling.  It carries none of the
Gibbs ringing at the limb that an FFT would impose, which is why the direct sum, and not an
FFT, is the reference for every certificate in this part.  The level radii come from a
bisection that brackets the \emph{inner} crossing of \eqref{eq:level-condition}: the left
side is not monotone, and the outer branch is the far side of the lobe.

\subsection{The principal axes and the ellipsoid}
\label{sec:deformation-axes}

Expanding the Legendre polynomials in \eqref{eq:level-depth} and simplifying puts the depth
in monomial form,
\begin{equation}
  E(n) \;=\; \tfrac32\Bigl[\epsilon_\text{tide}\,(n\cdot\hat e_\text{tide})^2
    - \epsilon_\text{spin}\,(n\cdot\hat e_\text{spin})^2\Bigr]
    + \tfrac{3\epsilon_\text{spin}-\epsilon_\text{tide}}{2} ,
  \label{eq:depth-form}
\end{equation}
every direction now entering as a squared projection.  Since
$(n\cdot\hat e)^2 = n^{T}(\hat e\hat e^{T})n$, the bracket is the quadratic form
$n^{T}Mn$ built from
\begin{equation}
  M \;=\; \tfrac32\Bigl(\epsilon_\text{tide}\,\hat e_\text{tide}\hat e_\text{tide}^{T}
    - \epsilon_\text{spin}\,\hat e_\text{spin}\hat e_\text{spin}^{T}\Bigr) ,
  \label{eq:depth-tensor}
\end{equation}
and the depth is
\begin{equation}
  E(n) \;=\; n^{T}Mn + \tfrac{3\epsilon_\text{spin}-\epsilon_\text{tide}}{2} .
  \label{eq:depth-quadratic}
\end{equation}
The constant does not spoil the identification.  On the unit sphere it equals $c\,n^{T}In$
with $c = \tfrac12(3\epsilon_\text{spin}-\epsilon_\text{tide})$, so
$E(n) = n^{T}(M + cI)\,n$: adding $cI$ shifts every eigenvalue by the same $c$ and moves no
eigenvector.  The principal directions are therefore the eigenvectors of $M$ itself, and the
constant only shifts the depths along them.
The principal axes of the level surface are those eigenvectors, and the radii along them
are the exact level radii from \eqref{eq:level-condition}.  Taking the ellipsoid through
those three points gives the shape the affine model of section~\ref{sec:deformation-stretch}
needs.

The temptation is to use the tide axis itself as the frame.  That is wrong whenever the
tide is weak or absent: with $\epsilon_\text{tide} = 0$ the tide's direction has no effect
on the potential, yet it would still rotate the frame, and on an oblate body the level
radius along a tilted direction differs from the equatorial one.  The shape would then
change with a parameter that cannot physically affect it.  The eigenvectors of
\eqref{eq:depth-tensor} do not have that defect, and when the equatorial eigenvalues are
degenerate the freedom they leave is harmless, because the radii there are equal.

\subsection{The stretch, and what it can and cannot model}
\label{sec:deformation-stretch}

The stretch replaces the map \eqref{eq:deformation-visibility} by a linear one.  Write the
deformed surface as $r(n)\,n$ with $n$ a unit direction and $n_\perp$ its sky-plane component;
the sky position of the emitting point is
\begin{equation}
  \mathbf{x}_\perp(n) \;=\; r(n)\, n_\perp .
  \label{eq:sky-map}
\end{equation}
That map is \emph{radial-nonlinear}: $r$ depends on the full direction, including its
component along the line of sight, so no constant matrix reproduces it.  What follows replaces
it by its linear part, $\Lambda n_\perp$.  For an ellipsoid the projection onto the sky is
exact and is an ellipse, so it is captured by a single $2\times2$ positive-definite matrix,
\begin{equation}
  \Lambda \;=\; \sqrt{A A^{T}} , \qquad A = P\,(\text{frame}\cdot\text{radii}) ,
  \label{eq:deformation-stretch}
\end{equation}
with $P$ the two rows $(\text{east},\text{north})$ projecting onto the sky and the radii
taken along the principal directions of \eqref{eq:depth-tensor}.  Writing $V_\text{sph}$ for
the visibility of a sphere carrying the same brightness pattern, the surrogate is
\begin{equation}
  V_\text{affine}(u,v) \;=\; V_\text{sph}\!\left(\Lambda^{T}(u,v)\right) .
  \label{eq:deformation-affine}
\end{equation}
There is \emph{no flux factor}.  The change of variables that stretches the image puts
$1/|\det\Lambda|$ on the brightness and $|\det\Lambda|$ on the measure, and the two cancel;
in the normalised quantity \eqref{eq:normalised-visibility} the determinant cancels
again.  So only the shape is modelled, while the brightness pattern --- gravity darkening
included --- is carried exactly.

$\Lambda$ is best read as a \emph{preconditioner} rather than as part of the answer.  It enters
twice --- in the frequency argument above and in the pull-back of the image --- and a different
but equally defensible $\Lambda$ changes the coefficients of any series built on it without
changing the exact answer; only the rate at which a truncated series converges depends on it.
Circularising the silhouette well is what makes a short series possible.

\paragraph{The same stretch in image space.}  The frequency change has a counterpart in the
image, and stating it makes the absent flux factor intuitive:
\begin{equation}
  V_\text{new}(u) = V(\Lambda^{T}u)
  \qquad\Longleftrightarrow\qquad
  I_\text{new}(y) = \frac{1}{|\det\Lambda|}\,I\bigl(\Lambda^{-T}y\bigr) .
  \label{eq:affine-image-transform}
\end{equation}
A feature at $x$ therefore appears at $y = \Lambda x$: the image is \emph{stretched} by
$\Lambda$, and its intensity is \emph{divided} by $|\det\Lambda|$, which exactly compensates
the area growth and leaves the total flux unchanged.  Since $\Lambda$ is symmetric positive
definite, $\Lambda^{T} = \Lambda$ and $\Lambda^{-T} = \Lambda^{-1}$: frequencies are
multiplied by $\Lambda$ while positions are divided by it.  For a general invertible $M$ the
two would be $Mu$ and $M^{-T}y$, which is why \eqref{eq:deformation-affine} is written with a
transpose; with this $\Lambda$ the transpose is cosmetic.

The model drops the Jacobian deliberately.  It moves the sample positions by $\Lambda$,
divides their brightness by $|\det\Lambda|$ and multiplies their weights by the same factor,
so the two cancel: what reaches the image is the stretched pattern with its brightness values
intact, and the image's total flux grows by $|\det\Lambda|$ instead of staying fixed.  That is
harmless because the discarded factor is a constant and $\Delta V$
\eqref{eq:normalised-visibility} is normalised at the origin --- and it is the honest
statement of the model's claim, which concerns shape and pattern rather than flux.

Two orderings have to be kept apart, and getting either wrong produces an error that looks
plausible.  The first is $\Lambda$'s own: it is expressed in the sky's
$(\text{east},\text{north})$ order, the order section~\ref{sec:conventions} gives the sky
axes and the order the image path uses, whereas the frequencies are $(u,v) =
(\text{North},\text{East})$.  The frequency is therefore reordered into the sky's order
before the stretch acts, and the result read back as $(\text{North},\text{East})$.  One
$2\times2$ matrix cannot serve both orders.  Using it in the wrong one transposes the
stretch, which is invisible for a nearly circular $\Lambda$ --- the pure-spin case --- and
wrong by order the anisotropy for anything else.

$V_\text{sph}$ is the mode sum of section~\ref{sec:mode-visibility}, evaluated on the
brightness pattern projected onto the harmonics.  Keeping $l \le 2$ is not a compromise: at
first order in the deformation the pattern is exactly $l = 0$ and $l = 2$, since the
$l = 4$ content first appears at second order from $P_{2,2}$.

No odd $l$ appears at any order, and the reason is symmetry rather than truncation.  Every
term of \eqref{eq:deformation-potential} is even under $n \to -n$ --- the monopole depends
on $r$ alone, and the others are squares of projections --- so $r(n)$ and $g(n)$ are even
functions of direction.  Since $Y_{lm}(-n) = (-1)^l Y_{lm}(n)$, an even function carries only
even $l$, and the argument is frame-independent because the point reflection is.  Spin and
the first-order tide therefore produce $l = 0$ and $l = 2$ and nothing else, in the body
frame and the observer frame alike; an $l = 1$ term would be a photocentre displacement,
which a deformation with this reflection symmetry cannot generate.  (A projection loop that
runs over $l = 1$ will still fill that slot, but with quadrature noise: it should shrink
with the sample count, and anything that does not is a real asymmetry and a bug.)

\paragraph{What a better shape would not fix.}  Expanding the level surface to the next order
\eqref{eq:level-condition} to second order in the deformation instead of stopping at the
ellipsoid, bringing $l = 4$ into the shape.  It would improve the outline at large amplitude,
where the ellipsoid visibly disagrees with the level surface.  It would \emph{not} repair the
stretch's error, because that error is first order and comes from the stretch, not from the
shape's second-order terms: a better ellipsoid does not make a linear map able to reproduce a
radial-nonlinear one.  Since the error is tolerable for the search the surrogate exists for, the second order
stays documented rather than implemented.

Should the surrogate itself need to improve, the productive direction is the basis rather than
the order in $\epsilon$: a series inside the same stretch would turn a fixed first-order error
into a truncation the user controls.

\subsection{The first-order row}
\label{sec:deformation-first-order}

\paragraph{The coefficients in closed form.}  To first order the pattern is a constant plus
two quadrupoles --- \eqref{eq:deformation-gravity-first-order} with $p = 1$ ---
\begin{equation}
  I \;=\; (1 - 4\epsilon_\text{spin})
    \;-\; 4\epsilon_\text{tide}\,P_2(n\cdot\hat e_\text{tide})
    \;+\; 4\epsilon_\text{spin}\,P_2(n\cdot\hat e_\text{spin}) ,
  \label{eq:first-order-pattern}
\end{equation}
Each $P_2(n\cdot\hat e)$ is a zonal $l = 2$ pattern about its own axis, so the corollary
\eqref{eq:zonal-l2} of section~\ref{sec:rotation-law} turns the bracket into harmonics
outright,
\begin{equation}
  I(\hat n) \;=\; 1 - 4\epsilon_\text{spin}
    - \frac{4}{5}\sum_{m=-2}^{2}\bigl(
      \epsilon_\text{tide}\,\overline{Y_{2m}(\hat e_\text{tide})}
      - \epsilon_\text{spin}\,\overline{Y_{2m}(\hat e_\text{spin})}\bigr)
      Y_{2m}(\hat n) ,
  \label{eq:first-order-harmonics}
\end{equation}
in the observer frame.  Nothing in that display is a coefficient yet: it is the pattern, with
the addition theorem applied and the terms left where they fell.  Collecting them, the
coefficients are
\begin{equation}
  a_{00} \;=\; 1 - 4\epsilon_\text{spin} , \qquad
  a_{2m} \;=\; -\frac{4}{5}\Bigl[\epsilon_\text{tide}\,
      \overline{Y_{2m}(\hat e_\text{tide})}
    \;-\; \epsilon_\text{spin}\,\overline{Y_{2m}(\hat e_\text{spin})}\Bigr] ,
  \label{eq:first-order-coefficients}
\end{equation}
every other coefficient being zero, with the two axes referred to the observer frame by the
same rule as a sampled direction, \eqref{eq:observer-angles}.  So this row of the model
needs \emph{no projection at all}: ten harmonic evaluations and one constant, with only the
level radii entering $\Lambda$ still numerical.  The first-order row of the model evaluates
\eqref{eq:first-order-coefficients} directly, so that row costs ten harmonic evaluations rather
than a projection.  Certificate \vtag{21} compares the production form against the quadrature
that projects the same pattern, so a wrong sign, constant or frame would show up as a residual
of order the coefficients rather than the $10^{-5}$ that discretisation leaves --- and because
the certificate calls the same function the model does, it covers the code rather than a copy
of it.

\paragraph{One formula, and a certificate for the frame.}  The pattern is fixed by two axes and
two amplitudes, all frame-independent objects, so \eqref{eq:first-order-coefficients} holds in
any frame: the model carries the axes as the body has them and rotates the coefficients
afterwards, while the row evaluates the same formula with the axes already rotated.  The two
must agree, and comparing them is \textbf{certificate \vtag{24}} --- which is what tests the
frame rule below rather than the formula itself, since both routes share the $l = 2$ corollary.

One is tempted to linearise the radii as well, replacing the bisection by $\delta r = -\phi_1$
and making the row analytic all the way through.  The temptation should be resisted.  First it
is exact only to first order, and the next term is not small: expanding
\eqref{eq:level-condition} gives
\begin{equation}
  \delta r \;=\; E + 3E^2 + \cdots ,
  \label{eq:radius-series}
\end{equation}
so at the amplitude ceiling, where $E = \epsilon_\text{tide} + 1.5\epsilon_\text{spin}$ reaches
$0.144$ against the pinch at $0.148$, the linearised radius is $0.144$ where the cubic gives
$0.366$ --- a bulge barely two fifths of the true one.  The series is asymptotic, and the
ceiling is precisely where it diverges.  Second, the saving would be three bisections per
parameter set, some fifty scalar evaluations, against the $10^5$ that the quadrature it
replaces used to cost: nothing measurable.  A shortcut that is safe only in a corner and buys
nothing is a hazard rather than an option, so it is not offered at all.  The model uses the
exact radii and linearises only the brightness, which keeps the difference between the two
analytic rows a genuine ``$O(\epsilon^2)$'' comparison rather than a shape error in disguise.

Each frequency therefore costs a
handful of Bessel functions, which is the whole point: the projection is paid once per
parameter set and the per-frequency cost is negligible, so a search over orientations is
cheap even at a sparse set of baselines.

The second ordering is the projection's, and it is worth stating because it is the point at
which the surrogate has most often gone wrong.  The pattern must be projected in the
\emph{observer} frame, not the body frame.  Let $(\text{north},\text{east})$ be the sky axes
of section~\ref{sec:conventions} for the viewing direction, and take as polar axis
\begin{equation}
  \hat n \;=\; \text{north} \times \text{east} \;=\; \hat e_\text{obs} ,
  \label{eq:observer-polar-axis}
\end{equation}
the direction \emph{toward} the observer.  That identity is not free: it holds only when the
sky's North is the projection of the body's North axis and the sky's East is
$\hat e_\text{obs}\times\text{North}$, the pairing section~\ref{sec:conventions} fixes.  Taking
East from a helper axis instead gives a North that is the projection of \emph{minus} North ---
the same handedness, but rotated by $\pi$ in the sky plane.  The consequence is a reflection of
the sphere, and since $Y_{lm}(\pi-\theta,\varphi) = (-1)^{l+m}Y_{lm}(\theta,\varphi)$ it leaves
every $|a_{lm}|$ unchanged while negating the $m = 1$ coefficients and leaving $m = 2$ alone:
the coefficients look plausible, every imaginary part has the wrong sign, and only a
diagnostic that compares two routes to the same number shows it.  Then
for each sampled direction $\mathbf{x}_i$,
\begin{equation}
  \theta_i = \arccos(\mathbf{x}_i\cdot\hat n) , \qquad
  \varphi_i = \operatorname{atan2}(\mathbf{x}_i\cdot\text{east},\, \mathbf{x}_i\cdot\text{north}) ,
  \label{eq:observer-angles}
\end{equation}
and $a_{lm} = \frac{1}{N}\sum_i I_i \overline{Y_{lm}(\theta_i,\varphi_i)}$ is a plain
quadrature, exactly as in section~\ref{sec:mode-visibility}.

Projecting in the body frame and rotating the coefficients afterwards does not do instead.
A rotation such as $R_y(i)$ fixes where $\hat z$ goes and leaves the azimuthal reference
somewhere else, so the frame is wrong by a rotation: the real and imaginary parts of $V$
then differ from the direct sum by a quarter turn.  In the axisymmetric case it fails more
quietly still --- the pattern stays axisymmetric about the line of sight and its visibility
is exactly real, so the imaginary part vanishes identically rather than being wrong.  Both
failures are what the undeformed null test and the direct-sum comparison of
section~\ref{sec:triaxial-register} are for.

The approximation is nonetheless \emph{first order}, and that is worth stating precisely
because it is easy to assume otherwise.  Comparing \eqref{eq:deformation-affine} against
\eqref{eq:deformation-visibility} on the same samples gives a residual of $0.0016$ at the
working amplitudes and $0.0208$ with the amplitudes halved, a ratio of $0.48$; with the
deformation switched off entirely the residual is zero to printed precision.  The error
therefore scales as $\epsilon$, not $\epsilon^2$, and it is the linearisation of the map
\eqref{eq:deformation-visibility} --- a stretch cannot reproduce a map that is
radial-nonlinear --- rather than a defect in the projection.

It is not, however, simply proportional to $\epsilon$.  At the amplitude ceiling, where
$E = 0.144$, the same comparison gives $0.3975$: the coefficient grows from about $1.3E$ to
$2.8E$ between the two regimes.  That is what an asymptotic series with a large next term
should do --- the radius expansion \eqref{eq:radius-series} already carries $3E^2$, and $E$
is approaching the pinch at $0.148$ where the surface opens.  So the surrogate is a
few-per-cent instrument for $\epsilon \lesssim 0.03$ and a forty-per-cent one at the ceiling.
For locating the interesting part of parameter space that is still usable, provided the
search is confined to the small-amplitude region; for deciding between models it is nowhere
near good enough, and the numerical model is what decides.

Written out, the pattern is multiplied by the limb factor:
\begin{equation}
  I(\hat n) \;=\; \bigl(1 - u_\text{limb}\,(1-\mu)\bigr)\,
  \Bigl[\,1 - 4\epsilon_\text{spin}
  - 4\epsilon_\text{tide}\,P_2(\hat n\cdot\hat t)
  + 4\epsilon_\text{spin}\,P_2(\hat n\cdot\hat s) \Bigr] ,
  \qquad \mu = \hat n\cdot\hat e_\text{obs} ,
  \label{eq:limb-pattern}
\end{equation}
whose coefficients are, to leading order in $u_\text{limb}\epsilon$,
\begin{equation}
  \begin{aligned}
    a_{00} &= (1-u_\text{limb})\,(1-4\epsilon_\text{spin}) , \\
    a_{10} &= \frac{u_\text{limb}}{\sqrt3} , \qquad
    a_{2m} = -\frac{4}{5}\Bigl[\,\epsilon_\text{tide}\,\overline{Y_{2m}(\hat t)}
    - \epsilon_\text{spin}\,\overline{Y_{2m}(\hat s)}\,\Bigr] .
  \end{aligned}
  \label{eq:limb-coefficients}
\end{equation}
None of these coefficients needs a Gaunt integral, and it is worth saying why, since the limb
factor is a product.  The dipole is the whole of the limb term because $\mu$ \emph{is} a single
harmonic: $Y_{10} = \sqrt3\cos\theta$, so $\mu = Y_{10}/\sqrt3$, and its
coefficient is its own amplitude.  The constant term is the other half of
$1 - u_\text{limb} + u_\text{limb}\mu$, multiplied by the pattern's mean.  The quadrupoles are
untouched, and come from the corollary \eqref{eq:zonal-l2} of
section~\ref{sec:rotation-law} rather than from any coupling of harmonics.

Coupling enters only through the terms that are dropped, and those are exactly the ones a Gaunt
coefficient describes: multiplying the factor into the pattern couples
$Y_{2m}\cdot Y_{10}$ into $l = 1$ and $l = 3$, with coefficients the $2\times1 \to 1$ and
$2\times1 \to 3$ integrals.  Two further corrections of order $u_\text{limb}\epsilon$ go with
them: the quadrupoles are not scaled by $(1-u_\text{limb})$.  What makes the omission defensible
is that the row is a model chosen for algebra and not a truncation of a physical law, and that
$a_{10}$ carries no deformation at all, so at zero tide and spin the row is exact --- which is
what \code{test\_at\_zero\_deformation\_the\_rows\_agree\_with\_the\_limb\_on} checks.

A quadratic law in $\rho^2$ was tried first and is too drastic: at half the disc radius it
gives $1 - 0.75u_\text{limb}$ where the linear law gives $1 - 0.5u_\text{limb}$.  Since the numeric row's law is a
placeholder for what a stellar-atmosphere model would supply, the linear form is the better
default: it is the classical one, and the closed-form row can carry it exactly at leading
order.  The factor multiplies the quadrupole pattern, so strictly the cross terms at
$O(\epsilon c)$ add small $l = 3$ structure; those are dropped, and the row stays a constant,
a dipole and two quadrupoles.

\subsection{The fitted row}
\label{sec:deformation-fitted}

\paragraph{The model in image space.}  The visibility is the transform of an image, and writing
that image explicitly is what makes the fit a projection.  By
\eqref{eq:affine-image-transform}, $V = \sum A_{lm}V_{lm}(\Lambda^{T}u)$ is the transform of
\begin{equation}
  I_\text{model}(x) \;=\; F\bigl(\Lambda^{-1}x\bigr) ,
  \qquad F(y) \;=\; \sum_{lm} A_{lm}\,\Phi_{lm}(y) ,
  \label{eq:fitted-model-image}
\end{equation}
with no $|\det\Lambda|$: the stretch is a coordinate change on the pattern, not a resampling
of a map, so the model is not required to conserve flux against another model star.  What is
wanted is that it match the star in front of it, and that is a fit, not a normalisation.
where the \textbf{disk harmonic} $\Phi_{lm}$ is the projection of $Y_{lm}$ onto the unit disk.  It
needs no integral and no Zernike expansion, because the projection \emph{is} the visible surface
value: the opaque convention absorbs the viewing cosine into the area element,
$d^{2}x = (\hat n\cdot\hat e)\,d\Omega$, so the image at a sky position is the brightness there, and
\begin{equation}
  \Phi_{lm}(x,y) \;=\; Y_{lm}(n) ,
  \qquad n \;=\; x\,\text{east} + y\,\text{north} + \sqrt{1-\rho^{2}}\;\hat e_\text{obs} ,
  \label{eq:disk-harmonic}
\end{equation}
with $Y_{lm}$ evaluated in the observer frame by the rule of \eqref{eq:observer-angles}.  There is no
cosine factor: for a uniform star $\Phi_{00}$ is constant, the image is the uniform disk whose flux
is $\pi$, and the visibility is the anchor of section~\ref{sec:visibility}.  The parametrisation
$\rho = \sin\theta$ makes every $\Phi_{lm}$ smooth on $\theta \in [0,\pi/2]$, the
$\sqrt{1-\rho^{2}}$ of the odd family being $\cos\theta$ there --- it comes from the parity of
$P_{lm}$ in its argument, not from a viewing cosine --- so inner products converge spectrally under
Gauss--Legendre.

\paragraph{The fit.}  Substituting $y = \Lambda^{-1}x$ in the image-plane residual gives
\begin{equation}
  \int \bigl|I_\text{num} - I_\text{model}\bigr|^{2}\,dx
  \;=\; \int_{|y|\le 1} \bigl|I'(y) - F(y)\bigr|^{2}\,|\det\Lambda|\,dy ,
  \qquad
  I'(y) \;=\; I_\text{num}\bigl(\Lambda y\bigr) ,
  \label{eq:fitted-unstretched}
\end{equation}
so the fit is a projection of the \emph{unstretched} image onto the disk harmonics in the sky
measure, and its normal equations are $G A = b$ with
\begin{equation}
  G_{lm,l'm'} \;=\; \int_{|y|\le 1} \Phi_{lm}\,\overline{\Phi_{l'm'}}\,dy ,
  \qquad
  b_{lm} \;=\; \int_{|y|\le 1} I'(y)\,\overline{\Phi_{lm}(y)}\,dy .
  \label{eq:fitted-normal-equations}
\end{equation}
$G$ carries no data: it is block diagonal in $m$ by the $e^{im\phi}$ factors, with blocks of size
$3$ for $m = 0$ ($l = 0, 2, 4$), $4$ for $m = 1, 2$ ($l = 2, 4$) and $2$ for $m = 3, 4$
($l = 4$) --- fifteen real parameters, matching $1 + 5 + 9$.

\paragraph{The right-hand side is a sum over the samples.}  Unstretching the data rather than the
basis,
\begin{equation}
  b_{lm} \;=\; \sum_s w_s\,I_s\,\overline{\Phi_{lm}\bigl(\Lambda^{-1}x_s\bigr)} ,
  \label{eq:fitted-rhs}
\end{equation}
so a fit costs one $2\times 2$ solve and fifteen mode evaluations per rendered sample.  No image
grid is built and nothing is interpolated.  The same sums give
\begin{equation}
  \bigl\|I\bigr\|^{2} \;=\; \sum_s w_s\,I_s^{2} ,
  \qquad
  \bigl\|I - M\bigr\|^{2} \;=\; \bigl\|I\bigr\|^{2} \;-\; b^{\dagger}G_s^{-1}b ,
  \label{eq:fitted-residual}
\end{equation}
with $M(x) = F(\Lambda^{-1}x)/|\det\Lambda|$ the model image and
$G_s = \sum_s w_s\,\Phi\Phi^{\dagger}$ built from the same samples as $b$.  The second is the
residual at the optimum, which is the number that makes this rung worth having.

$G_s$ is the sample Hessian rather than the disk Gram $G$ of
\eqref{eq:fitted-normal-equations}, and the
difference matters: $G$ is the basis's own integral over the unit disk, while the samples'
weights are projected areas, so the two agree only in the continuum limit.  Fitting with $G$
makes the reported residual optimal up to that measure rather than exactly, and adding a degree
can then \emph{raise} the number --- which is how the distinction was found.  $G$ remains what
the condition numbers describe, since those are a property of the basis alone.

\paragraph{Bookkeeping.}  The map is real, so $A_{l,-m} = \overline{A_{lm}}$ and only
$m \ge 0$ is fitted: one design function $\Phi_{l0}$ for $m = 0$, and the pair
$\mathrm{Re}\,\Phi_{lm}$, $\mathrm{Im}\,\Phi_{lm}$ for $m > 0$.  Odd $l$ is excluded, and not for
convenience: the parity argument of section~\ref{sec:deformation-brightness} forbids it, so fitting it
would only absorb noise.

\paragraph{What the row costs, and what it assumes.}  The projection is paid once per parameter
set --- $O(N(l+1)^2)$ for $N$ samples --- and each frequency then costs a short Bessel sum over
the $(l+1)^2$ coefficients, which is the resource that is cheap.  Nothing in the per-frequency
path depends on the sampling.

Two assumptions are worth stating, because they bound what the row can be trusted for.  The
model image is $F(\Lambda^{-1}x)/|\det\Lambda|$ on the unit disc, so the \emph{model's}
silhouette is exactly the stretched circle, while the data's is the true deformed silhouette,
which is not an ellipse; the fit absorbs the difference into the coefficients rather than
resolving it.  And $\Lambda$ comes from the three-point ellipsoid of
section~\ref{sec:deformation-axes}, so a different but equally defensible $\Lambda$ changes the
coefficients --- and because the fitted row's brightness is painted onto the surface, they are
meaningful only at the orientation they were fitted for.

The row is therefore already a disk series: $F(y) = \sum_{lm}A_{lm}\Phi_{lm}(y)$ with
$\Phi_{lm}$ the projected harmonics, composed with the same stretch.  A basis whose transforms
were cheaper --- a Zernike series, say --- would be a different choice of basis in the same
construction, not a different construction.

\paragraph{Checks built into the definition.}  The residual \eqref{eq:fitted-residual} is
evaluated by two routes --- through the normal equations and by expanding the quadratic directly
--- and the two must agree: they differ only through the identity $A = G_s^{-1}b$, so a stray
factor of $|\det\Lambda|$ would leave the coefficients plausible and show up there.
Certificate \vtag{26} removes the approximation altogether: the samples are manufactured from a
known coefficient set using the disk's own quadrature weights, so every step is an identity and the
fitted coefficients must come back at roundoff.  That pins the unstretching, the real pairing and
the normalisation at once, with no renderer and no sampling noise to hide behind.  Certificate
\vtag{27} pins the basis instead: $G_{00,00} = \pi$ exactly, the disk's area, since
$\Phi_{00} = 1$, entries
between different orders vanish by the orthogonality of the angular factors, and the whole matrix
agrees with a brute-force midpoint rule over the disk that shares nothing with the Gauss--Legendre
construction but the definition.

\paragraph{What the numbers mean.}  The fitted $A_{lm}$ are \emph{not} the pattern's coefficients of
section~\ref{sec:deformation-stretch}: they are the coefficients of the unstretched image in the disk
basis.  For a sphere the two coincide, which is exactly what the round trip tests.  With $\Lambda$
fixed, the residual is the truncation floor \emph{given} the ellipsoidal shape, so comparing it with
the affine row's error separates the two errors that row conflates --- the shape model's own and the
truncation at $l \le 2$.  It is not the floor of the class: that would need $\Lambda$ fitted too, and
$\Lambda$ is only weakly identifiable, its limit being independent of it.  Finally, the five condition numbers of $G$ are what to watch as the
mode count grows; they are properties of the basis alone, fixed once rather than per parameter set.

\subsection{How good the first-order row is}
\label{sec:first-order-residual}

Write $g_\text{lin} = 1 + 4\phi_1$ for the row's brightness before the limb factor and
$q = 4\phi_1 + 4\epsilon_\text{spin}$ for its quadrupole part, so that $g_\text{lin} = q +
(1-4\epsilon_\text{spin})$ splits into quadrupole and constant.  The row applies the limb factor
\eqref{eq:limb-factor} to the constant and the limb dipole and leaves the quadrupoles bare,
which is what \eqref{eq:limb-coefficients} says.  Its two errors are then
\begin{equation}
  I_\text{exact} - I_\text{row} \;=\;
    \bigl(1 - u_\text{limb}(1-\mu)\bigr)\bigl(g - g_\text{lin}\bigr)
    \;-\; u_\text{limb}(1-\mu)\,q ,
  \label{eq:row-error}
\end{equation}
both brightnesses: the first of order $\epsilon^2$, the linearisation of $g$; the second of order
$u_\text{limb}\epsilon$, the limb factor that the quadrupoles never receive.

The residual is quoted relative, in the fitted row's metric,
\begin{equation}
  \rho \;=\; \sqrt{\frac{\sum_s w_s \bigl(I_s - I^\text{row}_s\bigr)^2}{\sum_s w_s I_s^2}} ,
  \label{eq:relative-residual}
\end{equation}
where $w_s$ is the area sample $s$ projects onto the sky plane --- the weight the renderer
gives it, because the visibility integrates over that plane and not over the surface.  Thus
$\rho = 10^{-2}$ means the row is good to a per cent of the map.

The row's brightness must be matched to the sampled surface's weighted mean before the
difference is taken.  The rows are normalised at the origin (D22), so their absolute brightness
is a convention rather than a model, and the surface's mean brightness differs from an analytic
row's at first order in the deformation; without that match $\rho$ measures the flux mismatch,
which is first order and swamps the linearisation entirely.  The match is the freedom the fit
exercises when it chooses its own constant, so $\rho$ is then in the same norm as the fitted
row's --- and the fit, being optimal over a span that contains the first-order coefficients,
must have the smaller of the two.  Its two terms scale
differently in the amplitudes,
\begin{equation}
  \rho \;\simeq\; A\,\epsilon^2 \;+\; B\,u_\text{limb}\,\epsilon ,
  \label{eq:rho-scaling}
\end{equation}
with $A$ and $B$ independent of them.  Halving both amplitudes, and letting the field follow so
that the level surface relaxes with it, divides the first term by four and the second by two:
$\rho(\epsilon/2)/\rho(\epsilon)$ is $1/4$ at $u_\text{limb} = 0$ and tends to $1/2$ as the limb
term dominates.  At $\epsilon \approx 0.03$ and $u_\text{limb} = 0.6$ it dominates by an order of
magnitude, so the ratio is informative only with the limb off.

Two quantities are \emph{not} halved with the amplitudes: $u_\text{limb}$, which describes the
atmosphere rather than the deformation, and $\Lambda$, which is not a parameter but follows from
the level radii.  Holding either fixed changes the experiment.

The panel quotes $\rho$ for the first-order row and for the fitted ones in the same relative
form, so the two are comparable.  It quotes the halving ratio only when $u_\text{limb} = 0$, and
says \emph{n/a at $u > 0$} otherwise: the ratio is an order test, and with the limb on the order
it tests is not the one in \eqref{eq:row-error}.

\subsection{Register}
\label{sec:triaxial-register}

\begin{center}
\begin{tabular}{l l l}
\hline
tag & what is checked & kind \\
\hline
\vtag{19} & the first-order relations between $\delta r$, $\delta g$ and $\phi_1$ & algebraic \\
\hline
\vtag{20} & the surface really is level, and $\nabla\Phi$ matches a central difference & numerical \\
\hline
\vtag{21} & the closed-form first-order harmonics against the quadrature & numerical \\
\hline
\vtag{24} & the first-order coefficients in the body and observer frames & algebraic \\
\vtag{26} & the fit returns a synthetic image's own coefficients & algebraic \\
\vtag{27} & the Gram matrix: analytic entry, block structure, midpoint rule & numerical \\
\hline
\end{tabular}
\end{center}

Certificate \vtag{19} pins the two relations of the header of this section: on a level
surface $\delta r = -\phi_1/g_0$ and $\delta g = \partial_r\phi_1$, with the two contributions
to the gravity perturbation carrying the same sign, because the surface's displacement
lowers the gravity and the quadrupole's gradient raises it.  Certificate \vtag{20} is the
numerical anchor: it evaluates $\Phi$ on the returned surface and checks that the gradient
there agrees with a finite difference, so a wrong sign in \eqref{eq:deformation-potential}
would fail it immediately.

The strongest check in this part is not a tagged certificate but the null case: a zero
deformation must render as the uniform disk, with flux $\pi$ and the profile
$2J_1(x)/x$ of section~\ref{sec:uniform-disk}.  The oblate case then reproduces both the
equatorial bulge and the reduced equatorial gravity from the one field, which is the
statement that shape and gravity are not independent.
